\section{Results}

\begin{table*}[t]
  \centering
  \caption{Mean active-minus-random effects within each trajectory, averaged
  over optimization seeds 42 and 43.  Positive is favorable for PSNR, SSIM,
  and $F_5$; negative is favorable for LPIPS and AbsRel.  No trajectory or seed
  is excluded.}
  \label{tab:main-results}
  \begin{tabular}{lrrrrr}
\toprule
Trajectory & $\Delta$PSNR $\uparrow$ & $\Delta$SSIM $\uparrow$ & $\Delta$LPIPS $\downarrow$ & $\Delta$AbsRel $\downarrow$ & $\Delta F_{5}$ $\uparrow$ \\
\midrule
KT0 & +1.906 & +0.038 & -0.060 & -0.103 & +0.051 \\
KT1 & -0.233 & -0.010 & +0.005 & -0.022 & +0.014 \\
KT2 & +0.689 & +0.062 & -0.049 & +0.047 & -0.066 \\
KT3 & +0.187 & +0.016 & +0.012 & +0.059 & -0.005 \\
\midrule
Mean & +0.637 & +0.027 & -0.023 & -0.005 & -0.002 \\
\bottomrule
\end{tabular}

\end{table*}

\paragraph{Appearance improves, geometry does not.}
The policy improves mean PSNR by 0.637 dB and SSIM by 0.027 while reducing LPIPS
by 0.023 (Table~\ref{tab:main-results}).  The appearance effect is heterogeneous:
PSNR improves in 5/8 pairs and 3/4 trajectory means, and its trajectory-cluster
interval is $[-0.023,1.476]$ dB.  Raw depth AbsRel is nearly unchanged at
$-0.0047$ with a $[-0.0653,0.0531]$ interval.  Most importantly, mean 5 cm
surface F-score changes by $-0.0016$, is favorable in only 2/4 trajectory
means, and has interval $[-0.0459,0.0368]$.  KT0 exhibits aligned appearance
and geometry gains; KT2 gains 0.689 dB PSNR but loses 0.066 F-score, directly
demonstrating the image--geometry mismatch.

\begin{table*}[t]
  \centering
  \caption{Seed-42 grand means over four trajectories.  Coverage controls use
  the same seed set, target budget, test frames, optimizer seed, and training
  schedule as \method.  $F_5$ is surface F-score at 5 cm.}
  \label{tab:controls}
  \begin{tabular}{lrrrrr}
\toprule
Method & PSNR $\uparrow$ & SSIM $\uparrow$ & LPIPS $\downarrow$ & AbsRel $\downarrow$ & $F_{5}$ $\uparrow$ \\
\midrule
Random & 15.321 & 0.639 & 0.683 & 0.565 & 0.227 \\
Temporal coverage & 16.174 & 0.659 & 0.657 & 0.565 & 0.234 \\
Pose coverage & 16.270 & 0.665 & 0.642 & 0.524 & 0.242 \\
Depth-gradient + pose & 15.810 & 0.661 & 0.659 & 0.567 & 0.221 \\
\bottomrule
\end{tabular}

\end{table*}

\paragraph{Simple coverage is a stronger control.}
At seed 42, \method{} improves PSNR over random by 0.489 dB but reduces $F_5$
from 0.227 to 0.221 (Table~\ref{tab:controls}).  Temporal and pose coverage both
have higher PSNR and $F_5$ than \method; pose coverage is best on both
co-primary grand means.  The prespecified control-superiority claim therefore
fails.  This comparison matters because pose novelty already contributes 65\%
of the proposed score: the remaining depth-gradient term does not demonstrate
incremental geometric value over a direct coverage rule.

\begin{table*}[t]
  \centering
  \caption{Post-hoc error-ranking diagnostics for the target-free
  depth-gradient proxy.  Error targets are used only for evaluation.  AUSE is
  area under the sparsification-error curve.}
  \label{tab:diagnostics}
  \begin{tabular}{lrrrrrr}
\toprule
Trajectory & Seed & Views & $\rho$ & AUROC & AUPRC & AUSE $\downarrow$ \\
\midrule
KT0 & 42 & 165 & 0.219 & 0.782 & 0.488 & 0.346 \\
KT0 & 43 & 165 & 0.244 & 0.795 & 0.503 & 0.329 \\
KT1 & 42 & 170 & 0.100 & 0.705 & 0.315 & 0.433 \\
KT1 & 43 & 170 & 0.177 & 0.720 & 0.327 & 0.428 \\
KT2 & 42 & 165 & 0.079 & 0.678 & 0.331 & 0.359 \\
KT2 & 43 & 165 & 0.149 & 0.716 & 0.378 & 0.419 \\
KT3 & 42 & 165 & 0.091 & 0.662 & 0.321 & 0.377 \\
KT3 & 43 & 165 & 0.071 & 0.672 & 0.328 & 0.437 \\
\midrule
Mean & -- & -- & 0.141 & 0.716 & 0.374 & 0.391 \\
\bottomrule
\end{tabular}

\end{table*}

\paragraph{Error prediction succeeds.}
Spearman correlation is positive and AUROC exceeds 0.5 in every seed model
(Table~\ref{tab:diagnostics}); mean values are 0.141 and 0.716, respectively.
Thus the failed acquisition result cannot be summarized as an entirely
uninformative proxy.  Rather, prioritizing high rendered depth gradients tends
to identify difficult views without reliably identifying views whose addition
improves the held-out surface.  Difficulty, reducible error, redundancy, and
geometric coverage are distinct quantities.

\begin{table}[t]
  \centering
  \caption{Prespecified decisions.}
  \label{tab:gates}
  \begin{tabular}{lr}
\toprule
Rule & Outcome \\
\midrule
H1: positive mean PSNR and $F_5$ & Fail \\
H2: replication sign counts & Fail \\
H3: coverage-control superiority & Fail \\
H4: error-ranking diagnostic & Pass \\
LPIPS safety & Pass \\
Overall support & Fail \\
\bottomrule
\end{tabular}

\end{table}

\paragraph{Decision gates.}
Table~\ref{tab:gates} reports every frozen rule.  H1 fails because mean $F_5$
is not positive.  H2 fails because both co-primary outcomes are favorable in
only 5/8 pairs and $F_5$ in only 2/4 trajectory means.  H3 fails against the
coverage controls.  H4 passes, as does the LPIPS safety gate: mean LPIPS
improves, and the largest trajectory-mean regression is 0.012, below the
prespecified 0.02 limit.  Because H1 and H2 fail, overall support is false.

\paragraph{Protocol correction.}
Before any confirmatory model completed or outcome was inspected, initial
configuration logs exposed two incompatible defaults: pose auto-normalization
would invalidate metric surface evaluation, and validation frames were included
in nominal training budgets.  The batch was terminated during initialization,
the materializer and training configuration were corrected, every dataset was
rematerialized, and the preregistration was amended with the canceled run and
replacement commit before training restarted.
